% =====================================================================
%  Thème 1 — Conversions & cercle trigonométrique — NIVEAU MOYEN
%  Exercices
% =====================================================================
\documentclass[11pt,a4paper]{article}
\usepackage[utf8]{inputenc}
\usepackage[T1]{fontenc}
\usepackage{lmodern}
\usepackage[french]{babel}
\usepackage{amsmath,amssymb,mathtools}
\usepackage{icomma}
\usepackage{xcolor}
\usepackage{colortbl}
\usepackage{tikz}
\usepackage[most]{tcolorbox}
\usepackage{enumitem}
\usepackage{array}
\usepackage[a4paper,margin=2.2cm]{geometry}
\usetikzlibrary{arrows.meta,calc}

\definecolor{primary}{RGB}{214,51,108}
\definecolor{primarydark}{RGB}{140,20,64}

\DeclareMathOperator{\tg}{tg}
\DeclareMathOperator{\cotg}{cotg}
\newcommand{\degr}{^{\circ}}
\newcommand{\Z}{\mathbb{Z}}

\newtcolorbox{consigne}{enhanced,breakable,colback=primary!5,colframe=primary,
  boxrule=0.8pt,arc=2pt,left=3mm,right=3mm,top=2mm,bottom=2mm}

\newcounter{exo}
\newcommand{\exo}[1]{\medskip\refstepcounter{exo}%
  \noindent{\large\bfseries\color{primarydark}Exercice \theexo}%
  \ \textcolor{primary}{\textbullet}\ \textbf{#1}\par\nopagebreak\smallskip}

\setlength{\parindent}{0pt}
\pagestyle{empty}

\begin{document}

\begin{center}
{\LARGE\bfseries\color{primary}Thème 1 — Conversions \& cercle trigonométrique}\\[2pt]
{\color{primarydark}\textsc{Niveau Moyen — Exercices}}
\end{center}
\hrule height 1pt
\medskip

\begin{consigne}
\textbf{Objectif :} enchaîner \emph{quadrant $\to$ signe $\to$ angle de référence $\to$ valeur}.
On travaille désormais hors du premier quadrant, avec des angles négatifs ou supérieurs à un tour.
\textbf{Sans calculatrice}, valeurs exactes.
\end{consigne}

\exo{Valeurs remarquables hors du premier quadrant}
Donne la valeur exacte (pense au signe puis à l'angle de référence).
\[
\textbf{(a)}\ \cos 120\degr \qquad
\textbf{(b)}\ \sin 225\degr \qquad
\textbf{(c)}\ \tg 150\degr \qquad
\textbf{(d)}\ \cos 300\degr \qquad
\textbf{(e)}\ \sin\frac{4\pi}{3}
\]

\exo{Ramener dans un tour, puis évaluer}
Ramène d'abord l'angle dans $[0\degr;360\degr[$ (ou $[0;2\pi[$) en ajoutant/retranchant des tours,
puis donne la valeur exacte.
\[
\textbf{(a)}\ \cos 405\degr \qquad
\textbf{(b)}\ \sin(-30\degr) \qquad
\textbf{(c)}\ \tg 390\degr \qquad
\textbf{(d)}\ \cos\frac{7\pi}{2} \qquad
\textbf{(e)}\ \sin\frac{13\pi}{6}
\]

\exo{Lire un point sur le cercle}
Sur le cercle trigonométrique ci-dessous, le point $P$ correspond à un angle $\alpha$.
\begin{center}
\begin{tikzpicture}[scale=1.9,line cap=round,>=Stealth]
  \draw[primary,line width=1pt] (0,0) circle (1);
  \draw[->,black!70] (-1.3,0)--(1.35,0) node[right]{\scriptsize$\cos$};
  \draw[->,black!70] (0,-1.3)--(0,1.35) node[above]{\scriptsize$\sin$};
  \draw[primarydark,line width=1pt] (0,0)--(135:1);
  \fill[primarydark] (135:1) circle (0.03) node[above left]{\scriptsize$P$};
  \draw[black!60] (0:0.25) arc (0:135:0.25);
  \node[font=\scriptsize] at (80:0.4) {$\alpha$};
  \fill (0,0) circle (0.02);
\end{tikzpicture}
\end{center}
\textbf{(a)} Dans quel quadrant se trouve $P$ ? \quad
\textbf{(b)} Sachant que $\alpha$ est remarquable, donne sa valeur en degrés et en radians. \quad
\textbf{(c)} Donne alors les coordonnées exactes $(\cos\alpha\,;\,\sin\alpha)$ de $P$.

\exo{Signe imposé, quadrant à trouver}
Pour chaque condition, indique le(s) quadrant(s) possible(s) pour l'angle $\alpha$.
\[
\textbf{(a)}\ \cos\alpha>0\ \text{et}\ \sin\alpha<0 \qquad
\textbf{(b)}\ \tg\alpha<0 \qquad
\textbf{(c)}\ \sin\alpha>0\ \text{et}\ \cos\alpha<0
\]

\exo{Vrai ou Faux}
Justifie brièvement chaque réponse.
\begin{enumerate}[label=\textbf{(\alph*)},leftmargin=2em,itemsep=1pt]
\item $\dfrac{5\pi}{6}$ correspond à $150\degr$.
\item $\cos\dfrac{3\pi}{4}=\dfrac{\sqrt2}{2}$.
\item Un angle de $-45\degr$ est dans le même quadrant que $315\degr$.
\item $\tg 180\degr=0$.
\end{enumerate}

\end{document}
